\section{Introducción}
La Inteligencia Artificial Generativa y las arquitecturas de Retrieval-Augmented Generation (RAG) han adquirido una creciente relevancia en el desarrollo de aplicaciones de software, al permitir integrar modelos de lenguaje con componentes convencionales de recuperación, procesamiento y gestión de información. Esta integración introduce nuevas formas de interacción entre componentes de software que plantean desafíos para la calidad, mantenibilidad y evolución de las aplicaciones.

En este contexto, el presente anteproyecto, presentado ante el grupo de investigación SIGIIP, propone estudiar la presencia de antipatrones arquitectónicos en aplicaciones basadas en Inteligencia Artificial Generativa y RAG, con énfasis en su caracterización y detección mediante análisis estático. El documento presenta la selección y delimitación del tema, el planteamiento y formulación del problema, los objetivos, la justificación, el marco de referencia y las hipótesis de trabajo que orientan la investigación propuesta.